\scanpage{388}\mspage{127}

d'où $\mathrm{comp}_{\tau_0}=\mathrm{comp}_{\tau'_0}$.

Je recule devant la vraie difficulté -- celle de prouver l'associativité de cette loi de composition. Elle devient associative dans $\widehat{\mathrm{Hot}_0}$ -- mais \deleted{\unsure{pour que deux}} \inserted{pour que deux} flèches dans $\mathrm{Hot}$ soient égales (ici deux \deleted{de} \deleted{$C_\infty\times C_\infty\times C_\infty$}

\[
C_\infty(a,b)\times C_\infty(b,c)\times C_\infty(c,d)\rightrightarrows C_\infty(a,d)),
\]

suffit-il qu'il en soit ainsi de leurs images dans $\widehat{\mathrm{Hot}_0}$ -- en d'autres termes, le foncteur

\[
\mathrm{Hot}\unsure{-|\longrightarrow}\widehat{\mathrm{Hot}_0}
\]

est-il (en plus d'être conservatif) au moins fidèle (sinon pleinement fidèle) ?

Il y a aussi que cette loi de comp. est compatible avec celle des $H(a,b)$ de $\mathrm{Ind\,Hot}$, pour les flèches canoniques

\[
H(a,b)\longrightarrow C_\infty(a,b)\qquad\text{dans }\mathrm{Ind\,Hot}.
\]

Si on savait que c'est des \unsure{épis|épi}, et le restent par produits cartésiens, on aurait gagné.

Bien sûr, \deleted{\unsure{c'est}} on saura que c'est associatif, une fois qu'on saura que \inserted{la composition définie ici} est compatible avec la composition des chemins dans le cadre $1/2$ simplicial. Mais comment fait-on dans ce cadre, pour composer les chemins de longueur \deleted{\unsure{(et de t|et de t}} non bornée ? N'y a-t-il pas là exactement la même difficulté ? À moins de \unsure{s'app|s'en}

% archive sheet number: 387
